\subsection{A quotient norm bound from an orthogonal projection}
\label{archimedean-quotient-mass}

The unweighted Bessel estimate has a short proof by positivity.  Only
one integrable matrix coefficient is needed.  The resulting bound has
constant one for the original Haar measures.

Let $G$ be a second countable unimodular Lie group, $S$ a closed
unimodular subgroup, and $dg,ds,dq$ the measures in the Weil formula
on $G$, $S$, and $Q=S\backslash G$.  Let $\sigma$ be an irreducible
unitary representation on $H$, with formal degree $d>0$ for $dg$.
Inner products are conjugate linear in the first argument, and the
formal-degree convention is
\[
 \int_G\langle a,\sigma(g)b\rangle
       \overline{\langle c,\sigma(g)e\rangle}\,dg
       =d^{-1}\langle e,b\rangle\langle a,c\rangle.
\]
Let $\lambda:S\to U(W)$, with $W$ finite dimensional, and let
$B:H^\infty\to W$ be continuous and linear for the smooth-vector
topology, with $B\sigma(s)=\lambda(s)B$.

\begin{theorem}\label{original-quotient-mass-bound}
Suppose a nonzero smooth vector $v_0$ has an integrable diagonal
matrix coefficient.  Put
$D_0=\operatorname{span}\{\sigma(g)v_0:g\in G\}$, and suppose the
following operator integral is absolutely convergent for $u,v\in D_0$:
\begin{equation}\label{quotient-mass-period}
 \int_S\langle v,\sigma(s)u\rangle\lambda(s)^{-1}\,ds
                  =B(u)B(v)^*.
\end{equation}
Here $a b^*$ denotes the rank-one map $z\mapsto a\langle b,z\rangle$.
Then, for every $v\in H^\infty$,
\begin{equation}\label{quotient-mass-result}
       d\int_Q\|B\sigma(g)v\|^2\,dq(g)\le\|v\|^2.
\end{equation}
No nonvanishing condition on $B$ is required.
\end{theorem}

\begin{proof}
Every $v=\sum_i a_i\sigma(g_i)v_0\in D_0$ has an integrable
diagonal coefficient: it is the finite sum
\[
 \langle v,\sigma(t)v\rangle
 =\sum_{i,j}\overline{a_i}a_j
       \langle v_0,\sigma(g_i^{-1}tg_j)v_0\rangle.
\]
Fix a nonzero such $v$ and put
\[
       p(t)=\frac d{\|v\|^2}
                    \overline{\langle v,\sigma(t)v\rangle}.
\]
Schur orthogonality gives $p*p=p$, where
$(p*p)(x)=\int_Gp(t)p(t^{-1}x)\,dt$; explicitly, the unscaled
integral is
$d^{-1}\|v\|^2\overline{\langle v,\sigma(x)v\rangle}$.
Unitarity also gives $\overline{p(t^{-1})}=p(t)$.

Consider the Hilbert space of square-integrable sections represented
by functions $F:G\to W$ satisfying $F(sg)=\lambda(s)F(g)$.
Its norm is $\int_Q\|F(g)\|^2dq(g)$, and right translation is
$(R(t)F)(g)=F(gt)$.  The Weil quotient measure is right invariant,
so these are unitary operators.  Since $p\in L^1(G)$,
\[
                      P=\int_Gp(t)R(t)\,dt
\]
is a bounded self-adjoint idempotent, hence an orthogonal projection.

Write $f(g)=B\sigma(g)v$.  This is a continuous equivariant section.
For a bounded measurable section $F$ with compact support in $Q$,
the convolution integral defines a representative of $PF$ and is
absolutely convergent at every $g$, since it is bounded in norm by
$\|F\|_\infty\|p\|_1$.  Substituting $x=gt$
and applying the Weil formula therefore gives
\[
 (PF)(g)=\int_QK(g,h)F(h)\,dq(h),\qquad
 K(g,h)=\int_Sp(g^{-1}sh)\lambda(s)\,ds.
\]
Both $\sigma(g)v$ and $\sigma(h)v$ belong to $D_0$.
Taking the adjoint of \eqref{quotient-mass-period}, with
$u=\sigma(h)v$ and $v=\sigma(g)v$, yields exactly
\begin{equation}\label{quotient-mass-kernel}
                    K(g,h)=\frac d{\|v\|^2}f(g)f(h)^*.
\end{equation}
The inverse in the period and its absence in this kernel are thus
accounted for by the adjoint.

The quotient map has continuous local sections.  They make the
associated vector bundle and its measurable sections explicit;
a countable cover also gives measurable representatives.  In
particular, $\|f(g)\|$ is a continuous function on $Q$, bounded on
each compact set.  If $E\subset Q$ is compact, then
$F=f\,1_E$ is a bounded square-integrable section, although its
support in $G$ need not be compact.  Put
$A_E=\int_E\|f\|^2dq$.  Formula~\eqref{quotient-mass-kernel} gives
\[
       \frac d{\|v\|^2}A_E^2
           =\langle F,PF\rangle_{L^2(Q)}
           \le\|F\|_{L^2(Q)}^2=A_E.
\]
If $A_E=0$ the required bound is immediate; otherwise division by
$A_E$ gives $d A_E\le\|v\|^2$.  Compact exhaustion of $Q$ proves
\eqref{quotient-mass-result} for $v\in D_0$.  The vector $v=0$
is immediate.  At no point has the integrability of $f$ been assumed.

Finally $D_0$ is dense in $H^\infty$ for its smooth topology.
Here is the required argument.  Its Hilbert closure is all of $H$
by irreducibility.  For $\phi\in C_c^\infty(G)$, the map
$\sigma(\phi):H\to H^\infty$ is continuous: each derived seminorm
is bounded by an $L^1$ norm of the corresponding derivative of
$\phi$ times the Hilbert norm of the input.  For $u\in D_0$,
the vector $\sigma(\phi)u$ is a limit in the smooth topology of
finite Riemann sums of translates of $u$, and hence lies in the
smooth closure of $D_0$.  Hilbert approximation of an arbitrary
$u\in H$ now gives the same assertion for $\sigma(\phi)u$.
A smooth approximate identity then approximates every
$u\in H^\infty$ in that topology.  Since it is a Fr\'echet
topology, choose $v_j\in D_0$ converging to $v$ in it.
Continuity of $B\sigma(g)$ gives pointwise convergence of the
Bessel functions, and Fatou's lemma gives
\[
 d\int_Q\|B\sigma(g)v\|^2dq
   \le\liminf_j d\int_Q\|B\sigma(g)v_j\|^2dq
   \le\lim_j\|v_j\|^2=\|v\|^2.
\]
\end{proof}

\begin{corollary}\label{original-quotient-mass-isometry}
Under the same hypotheses, if $B$ is nonzero on $H^\infty$, then
\eqref{quotient-mass-result} is an equality for every smooth vector.
\end{corollary}
\begin{proof}
For nonzero $v\in D_0$, the function $f=B\sigma(\,\cdot\,)v$
cannot vanish identically: the smooth-density argument above,
applied to the orbit of $v$, would then force $B=0$.  Continuity
therefore supplies a compact $E\subset Q$ with $A_E>0$.
We already know that $f\in L^2(Q)$.  For $F=f1_E$ and
$\alpha=d/\|v\|^2$, the kernel formula gives $PF=\alpha A_Ef$.
The projection identity gives
\[
 \alpha^2A_E^2\|f\|_2^2=\|PF\|_2^2
      =\langle F,PF\rangle=\alpha A_E^2.
\]
Division by the positive number $\alpha A_E^2$ gives
$d\|f\|_2^2=\|v\|^2$.  Finally, smooth approximation $v_j\to v$
and the already proved upper bound applied to $v_j-v$ give
$\|B\sigma(\,\cdot\,)(v_j-v)\|_2\to0$.  Passing to the limit
proves equality for every smooth $v$, including $v=0$.
\end{proof}

\paragraph{Why the original gap supplies the integrable vector.}
For $G=U(k+r,r+1)$ put $n=2r+k+1$ and $q=r+1$.  The previously
proved lowest compact-type estimate \cite{quotient-mass-lowest-type}
provides a smooth unit vector $v_0$ with
\[
 |\langle v_0,\sigma(g)v_0\rangle|
 \le\exp\left[-\frac h8
          \min\left(t(g)^2,\frac{t(g)}{\sqrt{2q}}\right)\right],
\]
where $t(g)$ is the Euclidean norm of the logarithms of the singular
values and $h$ is the gap between parameters of opposite signature
signs.
Since $t(g)\ge\log\|g\|_{\rm op}$ and
$\min(t^2,t/c)\ge t/c-1/(4c^2)$ for $c>0$,
\begin{equation}\label{quotient-mass-decay}
 |\langle v_0,\sigma(g)v_0\rangle|
      \le e^{h/(64q)}\|g\|_{\rm op}^{-a},
             \qquad a=\frac h{8\sqrt{2q}}.
\end{equation}

For completeness, the needed volume bound is elementary.
In the real vector space $M_n(\mathbb C)$, of dimension $2n^2$,
take $U=\{x:\|x-I\|_{\rm op}<1/4\}$.  Left multiplication by
$g\in G$ preserves ambient Lebesgue volume, because its real
Jacobian is $|\det g|^{2n}=1$.  Thus disjoint sets $g_iU$ with
$\|g_i\|_{\rm op}\le R$ all have the same positive volume and lie
in the ambient ball of radius $5R/4$.  There are at most
$C R^{2n^2}$ of them.  A finite maximal disjoint family covers
the group norm ball by the sets
$g_i(\overline{UU^{-1}}\cap G)$.  The factor in parentheses is
compact, since $U$ is bounded and bounded away from singular
matrices.  Haar measure of the group norm ball is consequently
at most $C'R^{2n^2}$.  Summing over dyadic shells proves
integrability of $\|g\|_{\rm op}^{-a}$ whenever $a>2n^2$.
This argument uses equal-volume ellipsoids, rather than a claim
that their Euclidean radii are fixed.

The original gap assumptions include
\[
 a\ge9(r+1)n^2+3A_G+3,
               \qquad A_G=4(k+r)(r+1),
\]
so $a>2n^2$.  Thus \eqref{quotient-mass-decay} supplies the needed
$L^1$ coefficient without an additional hypothesis or a uniform
bound for its $L^1$ norm.  That auxiliary constant disappears when
$P$ is recognized as an orthogonal projection.
For the original family of integer shifts, the supplied gap at a
shifted parameter is $h/4$ and the supplied numerical hypothesis is
the full gap condition at $h/4$.  Apply the preceding argument with
that value.  This retains the same shift reserve; it does not use
the larger unshifted gap in its place.

\paragraph{Original Bessel normalization and quotient measure.}
In the original setting $S=U(W_0)N$ and
$\lambda(hn)=\psi(n)\pi_W(h)$, with the given unitary character
and compact representation.  The subgroup is closed and unimodular.
Its Haar measure is the pushforward of the original product measure
on $U(W_0)\times N$.  The supplied subgroup-period normalization
is \eqref{quotient-mass-period}; its second input is allowed in the
span of translates of compact-finite vectors, which contains $D_0$.
The original finite sum of derived seminorms bounding $B$ gives
precisely the smooth continuity used above
\cite{quotient-mass-original-normalization}.

If the source writes the quotient integral using a measure $\nu$
on matrix representatives, the same conclusion holds for that
measure.  Its Weil identity first forces it to be supported on
$G$.  Its pushforward to $Q$ is $dq$: for a compactly supported
function $\varphi$ on $Q$, choose $b\in C_c(G)$, $b\ge0$, with
$A b(g)=\int_S b(sg)ds>0$ on its support, using finitely many local
sections.  Apply the two Weil formulas to
$b(g)\varphi(Sg)/A b(g)$.  This proves equality on compactly
supported test functions, local finiteness, and hence equality
of the quotient measures.  All nonnegative invariant integrals
therefore agree.  No Haar measure has been renormalized.

In particular $F_v=\sqrt d\,B\sigma(\,\cdot\,)v$ satisfies the
unweighted estimate used in the flag moments and the lowering
recursion.  When $B\ne0$, the corollary also gives the exact
isometry required by the original normalization.  Existence
of the original discrete series, the classical lowest compact-type
formula, and the supplied subgroup-period Bessel map remain
representation-theoretic inputs.  No weighted orbital estimate or
leading-minor identification is asserted here.
